%!TEX root = tfg.tex
\chapter{Introducción}

\begin{quotation}[Polymath]{Benjamin Franklin (1706--1790)}
Well done is better than well said.
\end{quotation}

\begin{abstract}
En esta parte se documenta la ejecución técnica del proyecto y su seguimiento a través
de iteraciones. Cada iteración recoge las funcionalidades trabajadas, las decisiones de
diseño, la implementación realizada y los resultados obtenidos.
\end{abstract}

\section{Enfoque iterativo}
La fase de desarrollo se ha organizado de forma incremental, siguiendo el flujo de trabajo
definido en la metodología del proyecto. Las funcionalidades del \textit{Feature List} se
han distribuido en iteraciones cerradas para mantener control sobre el avance, revisar la
integración entre sistemas y disponer de entregas jugables verificables.

\section{Seguimiento y control}
El seguimiento del progreso se ha realizado a través del \textit{backlog}, del tablero
Kanban y de las revisiones del repositorio. Esta combinación ha permitido relacionar las
decisiones de planificación con los cambios técnicos concretos y justificar el cierre o la
división de funcionalidades cuando ha sido necesario.
